\begin{theorem}\label{night:P3F.20}
\StatementRecord{night:P3F.20}{proved}{}{\cref{night:P3F.1}, \cref{night:P3F.21}, \SeriesRef{night:P3F.31}, \cref{night:P3K.65}; \cref{iv:rank_two_specialization}; \cite[Theorem 1.1, p.~2; Propositions 2.7--2.8, pp.~4--6; Sections 4--6, pp.~9--15]{Wang2026}}
For Wang's threefold, the primitive \(I_2\) filling has the
literal coordinate and specialization maps below, and its physical
line circle has isotropy \(\mathbb Z/2\) on the pinching locus.
The global conclusions \(\pi_1=0\) and
\(H_2=H_3=\mathbb Z/2\) are
\cite[Theorem 1.1, p.~2; proof in Sections 4--6, pp.~9--15]{Wang2026}.
The comparison retains the specified primitive filling and the
component-supported OS49 application stated in the proof.
\end{theorem}
\begin{proof}
\begin{enumerate}[label=\arabic*.]
\item Wang's cited construction uses source configuration
\(IV^*,I_2,I_1,I_1\), line \(M=\mathcal O(P-O)\), translation
by \(-6P\), and one order-three logarithmic cap.  Its topology is
the cited theorem.  Its nonzero \(H_2\) obstructs a homeomorphism
to \(S^6\) by the induced map on integral homology.
\item Retain the actual local periods
\[
h(z,v)=(q^2c_1z,qc_2v),\qquad
g(z,v)=(c_3z,\lambda c_4v),
\]
where the \(c_i\) are holomorphic units and the second generator
is the original lift multiplied by \(h^3\).  The invertible
monomial map
\[
 (z,v)\longmapsto(x,y)=(v,v^2/z),\qquad
 (x,y)\longmapsto(z,v)=(x^2/y,x)
\]
sends the exponent \((2,1)\) to \((1,0)\).  The zero-exponent
period multiplies \(y\) by \(\lambda^2c_4^2/c_3\), retaining
its strict contraction.  The ray fan therefore gives the actual
one-component primitive rank-one collapse.  Its surviving one-cell
and its two-cell both have coefficient one.  The other rank-one
filling has the same primitive collapse; the unit-bearing rank-two
filling has the full specialization morphism of
\cref{iv:rank_two_specialization}.  This supplies the explicit
local comparison required here, including the integral indices.
\item On the angular two-torus, the physical line circle is
\(\theta\mapsto(0,\theta)\), and the pinched circle is
\(a\mapsto(2a,a)\), both modulo \(\mathbb Z^2\).
Their intersection is obtained from \(2a\in\mathbb Z\):
\(a=0,1/2\) modulo one.  The physical line circle therefore has
exact stabilizer \(\{0,1/2\}\cong\mathbb Z/2\) on that
pinching locus.  This follows from the actual primitive collapse,
so the order-two isotropy is retained in further orbit comparisons.
\item For the component-supported OS49 source with \(Q=3P\),
the regular-binding degree-three page has three \(II\) and three
\(I_2\) fibres.  Thus \(\chi(\mathcal O)=1\), \(\sigma=-8\),
\(h(P)=1/2\), and the reducible-fibre correction is \(3/2\).
The height formula gives \(P\cdot O=0\).  Its actual
component-permuting translation and positive-fan page are carried
through \SeriesRef{night:P3F.31}, \SeriesRef{night:P3F.32}, giving index zero.
The complete native closing and ruling map in \cref{night:P3K.65}
gives negative-fan class \(-3\) for \(\lambda=r^3\) and
\(-1\) for the constructed shorter loop \(\lambda=w\).
These are the specified maps and page data; the source-configuration
enumeration in \cref{night:P3F.1} supplies no numerical bounds
for an additional bounded slice of all compact fillings.
\end{enumerate}
\end{proof}




